%%%%%%%%%%%%%%%%%%%%%%%%%%%%%%%%%%%%%%%%%%%%%%%%%%%%%%%%%%%%%%%%%%%%%%%%%%%%%%%%
% Theory of Quantum Mechanics
% Chapter 1 — Mathematical Language of Physics
% Section 1.1 — Mathematical Language of Physics
% Version 1.0 — integrated manuscript
%%%%%%%%%%%%%%%%%%%%%%%%%%%%%%%%%%%%%%%%%%%%%%%%%%%%%%%%%%%%%%%%%%%%%%%%%%%%%%%%

\section{Mathematical Language of Physics}
\label{sec:mathematical-language-of-physics}

Physics seeks to answer some of the most fundamental questions about Nature.
Although these questions arise from observation, their answers cannot be
expressed with sufficient precision in ordinary language alone. The remarkable
success of modern physics stems from the discovery that the laws governing the
Universe can be formulated as mathematical relationships whose predictions can
be tested experimentally.

Mathematics is therefore far more than a computational tool. It provides the
conceptual language in which physical theories are formulated, analyzed,
compared, and extended. In this chapter we begin developing that language,
starting with measurable quantities and geometric vectors and then moving
toward the abstract structures required by quantum mechanics.

\subsection{Why Mathematics?}
\label{subsec:why-mathematics}

Long before modern science, humanity recognized regular patterns in Nature: the
alternation of day and night, the changing seasons, the motion of celestial
bodies, and the fall of terrestrial objects. Such observations made calendars,
navigation, and the prediction of eclipses possible. Qualitative descriptions,
however, could not by themselves reveal the precise relationships connecting
these phenomena or generate reliable predictions in unfamiliar situations.

The emergence of theoretical physics required measurement, experiment, and
mathematical reasoning to become inseparable. Galileo showed that motion could
be studied quantitatively. Newton expressed terrestrial and celestial mechanics
through the same dynamical principles. Maxwell united electricity, magnetism,
and optics in a system of field equations. Einstein reformulated the concepts
of space, time, and gravitation. Schr\"odinger and Heisenberg introduced new
mathematical structures capable of describing microscopic phenomena for which
classical trajectories were no longer adequate.

In each case, mathematics did more than summarize observations. Once a physical
principle had been formulated mathematically, consequences could be derived
that were not evident from the original experiments. A theory could therefore
predict new phenomena, expose hidden relations among apparently different
processes, and identify the limits within which an older description remained
valid.

Mathematics is universal in a further sense. An equation does not depend on the
spoken language of the scientist who writes it. Its symbols acquire meaning
through definitions and operational rules, allowing the same physical law to be
examined, reproduced, and criticized across cultures and generations.

Experiment nevertheless remains the final judge of a physical theory.
Mathematics contains many internally consistent structures, but only a small
subset describes the physical world. Theoretical physics therefore advances
through a continuing dialogue:

\begin{center}
observation $\longrightarrow$ model $\longrightarrow$ prediction
$\longrightarrow$ experiment $\longrightarrow$ revision.
\end{center}

\begin{bookinsight}[Mathematics and physical reality]
A mathematical theory of physics is neither a collection of measurements nor a
free construction detached from experiment. It is a structured model whose
concepts are tied to possible observations and whose consequences must survive
experimental tests.
\end{bookinsight}

The progression of this chapter reflects that viewpoint. We begin with the
quantities obtained through measurement, pass through geometry and vectors,
and gradually move toward vector spaces, matrices, eigenvalue problems, inner
products, and operators. The final destination is the mathematical language in
which a quantum state can be understood as a vector and a measurable quantity
as an operator.

\subsection{Physical Quantities}
\label{subsec:physical-quantities}

Every physical theory begins with quantities that can, at least in principle,
be connected to measurement. Length, duration, mass, electric charge,
temperature, momentum, and energy are familiar examples. To specify such a
quantity completely, a numerical value is generally not enough; one must also
state the unit and, for directional quantities, the relevant direction or
components.

Thus the statement
\[
L=2.0\ \mathrm{m}
\]
contains both a number and a unit. The number tells us how many copies of the
chosen standard are present, while the unit identifies the standard against
which the comparison has been made. Changing the unit changes the numerical
value but not the underlying physical length:
\[
2.0\ \mathrm{m}=200\ \mathrm{cm}.
\]

Physical quantities may be classified as \emph{fundamental} or \emph{derived}.
A chosen system of units begins with a small set of base quantities. Other
quantities are then defined through physical or geometric relations. Velocity,
for example, is constructed from displacement and time, while force is related
to mass and acceleration.

\begin{bookdefinition}[Physical quantity]
A physical quantity is a property of a system or process that can be represented
quantitatively through a specified measurement procedure, together with an
appropriate unit and mathematical type.
\end{bookdefinition}

The phrase \emph{mathematical type} is important. Temperature and energy can be
represented by scalars, whereas displacement and momentum require vectors.
Later, quantum states will require a still more general kind of vector, and
observables will be represented by operators. The mathematical object chosen to
represent a quantity must preserve the structure observed in the physical
phenomenon.

\subsection{Measurements and Units}
\label{subsec:measurements-and-units}

Measurement is a comparison between a physical system and an agreed standard.
No real measurement produces an infinitely precise number. Every result is
limited by the instrument, the procedure, the surrounding conditions, and the
statistical variability of repeated observations. A measured quantity should
therefore be understood not merely as a number, but as a reported value together
with information about its uncertainty.

A simple measurement may be written schematically as
\[
q=q_{\mathrm{meas}}\pm\Delta q,
\]
where $q_{\mathrm{meas}}$ is the reported value and $\Delta q$ characterizes the
measurement uncertainty. The interpretation of $\Delta q$ depends on the
experimental and statistical context; a careful theory of uncertainty will be
developed later in the book.

Units make comparisons reproducible. The International System of Units, or SI,
uses seven base units, including the metre, kilogram, second, ampere, kelvin,
mole, and candela. Derived units abbreviate recurring combinations of these
base units. For example,
\[
1\ \mathrm{N}=1\ \mathrm{kg\,m\,s^{-2}},
\]
so the newton is a convenient name for the SI unit of force.

Units also provide an immediate consistency test. An equation intended to
represent a physical law must be dimensionally homogeneous: quantities added to
one another must have the same physical dimensions, and the dimensions on both
sides of an equality must agree.

\begin{bookwarning}[A number without its unit]
Except for dimensionless quantities, a numerical value written without its unit
is incomplete. The statements $L=2$, $t=5$, and $E=10$ do not yet specify
physical quantities.
\end{bookwarning}

The next section introduces the first major distinction among mathematical
representations of physical quantities: the distinction between scalars and
vectors.
